\begin{frame}{Stav práce}{Popis progresu}
\textbf{Fáze:} Testování distribuovaných úložišť \\
\vspace{0.5cm}
\begin{columns}[T]
\begin{column}{0.48\textwidth}
 \textbf{Testované OS:}
 \begin{itemize}
  \item CentOS Stream 
  \item Proxmox
  \item Rocky Linux 
  \item Ubuntu
  \item Windows 
 \end{itemize}
\end{column}
\begin{column}{0.48\textwidth}
 \textbf{Testované FS:}
 \begin{itemize}
  \item BeeGFS
  \item Ceph
  \item GlusterFS 
  \item Lustre
  \item NFS + Pacemaker
 \end{itemize}
\end{column}
\end{columns}
% ===============================
% Notes
% ===============================
\note{
\textbf{Testované OS:}
\begin{itemize}
  \item CentOS Stream (9, 10)
  \item Proxmox (8.2, 8.3, 9.0) 
  \item Rocky Linux (9.5, 9.7, 10.0)
  \item Ubuntu (22.04)
  \item Windows 
\end{itemize}

\begin{itemize}
	\item začátek na Proxmox(nemá podporu pro IPoIB a RDMA)
	\item Ubuntu(max 2 vlánka - 30Gb/s)
	\item CentOS(max 4 vlánka - 60Gb/s) (Bohuzel nejde o serverovy OS -> moznost nestability)
	\item Rocky Linux(max 4 vlánka - 89Gb/s) (podporuje RDMA a IPoIB) -> 9.5, 9.7, 10.0
\end{itemize}
}
\end{frame}


